% Source: https://artofproblemsolving.com/community/c3h1357872_reduction_of_sillies
% Stable identity: aops-topic-1357872-q01
Go through the list.
\begin{itemize}
\item Rational: $-2$, $0$, $-0.1111\ldots=-\dfrac{1}{9}$, $3.7256=\dfrac{37256}{10000}$, $\sqrt9=3$, $-\dfrac{3}{124}$.
\item Irrational: $\pi$; $\sqrt3$ ($3$ is not a perfect square); $0.112123123412345\ldots$ (its decimal expansion never becomes periodic, since the blocks keep growing); $\sqrt{2\pi}$ (if it were rational, $\pi=\dfrac{(\sqrt{2\pi})^2}{2}$ would be rational).
\end{itemize}
So $\boxed{6}$ of the numbers are rational. A good habit against slips like $\sqrt9$: simplify every expression first, then classify.
